\documentclass[12pt]{article}
\usepackage{amsmath, amssymb, amsthm}
\usepackage{geometry}
\usepackage{graphicx}
\usepackage{enumitem}
\geometry{margin=1in}

\title{CSE400 --- Fundamentals of Probability in Computing\\
Lecture 11--12: Transformation of Random Variables}
\author{Instructor: Dhaval Patel, PhD}
\date{}

\begin{document}

\maketitle

\noindent
\textbf{Course Context:} Transformation techniques for single and multiple random variables\\
\textbf{Assumption (as per slides):} PDF/CDF of original random variable(s) is known \textit{a priori}

\section*{1. Transformation of a Single Random Variable}

We begin with the fundamental problem:

\begin{quote}
Given a random variable $X$ with known distribution, and a new random variable defined as
\[
Y = g(X),
\]
find the distribution (CDF and PDF) of $Y$.
\end{quote}

\subsection*{1.1 CDF-Based Method (As Derived in Slides)}

We start from first principles.

\subsubsection*{Step 1: Write the CDF of $Y$}

\[
F_Y(y) = P(Y \le y)
\]

Since $Y = g(X)$,

\[
F_Y(y) = P(g(X) \le y)
\]

At this stage, the derivation depends on the monotonicity of $g(\cdot)$.

\subsection*{Case 1: $g(\cdot)$ is Monotonically Increasing}

If $g$ is strictly increasing, then

\[
g(X) \le y \quad \Longleftrightarrow \quad X \le g^{-1}(y)
\]

Therefore,

\[
F_Y(y) = P(X \le g^{-1}(y)) = F_X(g^{-1}(y))
\]

This is the key CDF transformation formula.

\subsubsection*{Step 2: Differentiate to Obtain the PDF}

\[
f_Y(y) = \frac{d}{dy} F_Y(y)
= \frac{d}{dy} F_X(g^{-1}(y))
\]

Using the chain rule:

\[
f_Y(y)
= f_X(g^{-1}(y)) \cdot \frac{d}{dy} g^{-1}(y)
\]

Thus:

\[
\boxed{
f_Y(y) = f_X(g^{-1}(y)) 
\left| \frac{d}{dy} g^{-1}(y) \right|
}
\]

(The absolute value accounts for decreasing transformations.)

\subsection*{Case 2: $g(\cdot)$ is Monotonically Decreasing}

If $g$ is decreasing:

\[
g(X) \le y \Longleftrightarrow X \ge g^{-1}(y)
\]

Thus,

\[
F_Y(y) = P(X \ge g^{-1}(y))
= 1 - F_X(g^{-1}(y))
\]

Differentiating:

\[
f_Y(y)
= - f_X(g^{-1}(y)) \frac{d}{dy} g^{-1}(y)
\]

Since derivative is negative, we again obtain:

\[
\boxed{
f_Y(y) = f_X(g^{-1}(y)) 
\left| \frac{d}{dy} g^{-1}(y) \right|
}
\]

\section*{2. Example from Slides}

\subsection*{Example: $X \sim \text{Uniform}(-1,1)$}

Find the PDF of:

\[
Y = \sin\left(\frac{\pi}{2}X\right)
\]

\subsubsection*{Step 1: Identify Transformation}

\[
y = \sin\left(\frac{\pi}{2} x \right)
\]

Since sine is increasing over $x \in (-1,1)$, we can invert.

\subsubsection*{Step 2: Invert}

\[
x = \frac{2}{\pi} \sin^{-1}(y)
\]

\subsubsection*{Step 3: Compute Derivative}

\[
\frac{dx}{dy}
=
\frac{2}{\pi} \cdot \frac{1}{\sqrt{1-y^2}}
\]

\subsubsection*{Step 4: Substitute into Formula}

Since

\[
f_X(x) = \frac{1}{2}, \quad -1 < x < 1
\]

we obtain

\[
f_Y(y)
=
\frac{1}{2}
\cdot
\frac{2}{\pi}
\cdot
\frac{1}{\sqrt{1-y^2}}
\]

Thus,

\[
\boxed{
f_Y(y) = \frac{1}{\pi \sqrt{1-y^2}}, \quad -1 < y < 1
}
\]

This matches the derivation shown in the lecture.

\section*{3. Function of Two Random Variables}

We now move to joint transformations.

Given:

\[
Z = X + Y
\]

Find the PDF of $Z$.

\subsection*{3.1 CDF-Based Derivation (As Presented)}

Start with:

\[
F_Z(z) = P(Z \le z)
\]

Since $Z = X + Y$,

\[
F_Z(z) = P(X + Y \le z)
\]

This defines a region in the $(x,y)$-plane:

\[
\{(x,y) : x+y \le z\}
\]

\subsubsection*{Step 1: Write as Double Integral}

\[
F_Z(z)
=
\iint_{x+y \le z}
f_{X,Y}(x,y) \, dx \, dy
\]

Region conversion:

\[
\int_{-\infty}^{\infty}
\int_{-\infty}^{z-x}
f_{X,Y}(x,y)
\, dy \, dx
\]

\subsubsection*{Step 2: Differentiate w.r.t. $z$}

\[
f_Z(z) = \frac{d}{dz} F_Z(z)
\]

Using \textbf{Leibniz Rule}:

If

\[
G(z) = \int_{a(z)}^{b(z)} h(x,z) \, dx
\]

then

\[
\frac{d}{dz}G(z)
=
h(b(z),z)b'(z)
-
h(a(z),z)a'(z)
+
\int_{a(z)}^{b(z)}
\frac{\partial}{\partial z} h(x,z) \, dx
\]

In our case:

\begin{itemize}
\item Lower limit does not depend on $z$
\item Upper limit is $z-x$
\end{itemize}

Thus,

\[
f_Z(z)
=
\int_{-\infty}^{\infty}
f_{X,Y}(x, z-x) \, dx
\]

\subsection*{3.2 Convolution Result (Independence Case)}

If $X$ and $Y$ are independent:

\[
f_{X,Y}(x,y) = f_X(x) f_Y(y)
\]

Therefore:

\[
\boxed{
f_Z(z)
=
\int_{-\infty}^{\infty}
f_X(x) f_Y(z-x) \, dx
}
\]

This is the convolution integral.

\section*{4. Example: Sum of Two Normal Random Variables}

Given:

\[
X \sim N(0,\sigma^2),
\quad
Y \sim N(0,\sigma^2)
\]

Independent.

\[
f_Z(z)
=
\int_{-\infty}^{\infty}
\frac{1}{\sqrt{2\pi}\sigma}
e^{-x^2/(2\sigma^2)}
\cdot
\frac{1}{\sqrt{2\pi}\sigma}
e^{-(z-x)^2/(2\sigma^2)}
dx
\]

Combining exponents:

\[
e^{-\frac{x^2+(z-x)^2}{2\sigma^2}}
\]

Final result:

\[
\boxed{
Z \sim N(0, 2\sigma^2)
}
\]

\section*{5. Example: Sum of Exponential Random Variables}

Given:

\[
X, Y \sim \text{Exponential}(\lambda)
\]

Independent.

\[
f_X(x) = \lambda e^{-\lambda x}, \quad x>0
\]

Using convolution:

\[
f_Z(z)
=
\int_0^z
\lambda e^{-\lambda x}
\lambda e^{-\lambda(z-x)}
dx
\]

\[
=
\lambda^2 e^{-\lambda z}
\int_0^z dx
=
\lambda^2 z e^{-\lambda z}
\]

Thus,

\[
\boxed{
f_Z(z)
=
\lambda^2 z e^{-\lambda z}, \quad z>0
}
\]

\section*{6. Example: Difference of Two RVs}

For:

\[
Z = X - Y
\]

\[
F_Z(z) = P(X - Y \le z)
\]

Region:

\[
x - y \le z
\]

Follow identical region-based integration approach.

\section*{7. Ratio Example (Mentioned in Slides)}

\[
Z = \frac{X}{Y}
\]

Also:

\[
R = \sqrt{X^2 + Y^2}
\]

These require Jacobian-based joint transformations.

\section*{Final Exam-Ready Summary}

\subsection*{1. Single RV Transformation}

\[
f_Y(y)
=
f_X(g^{-1}(y))
\left|
\frac{d}{dy}g^{-1}(y)
\right|
\]

\subsection*{2. Sum of Two RVs}

\[
F_Z(z)
=
P(X+Y \le z)
\]

\[
f_Z(z)
=
\int f_{X,Y}(x,z-x) \, dx
\]

If independent:

\[
f_Z(z)
=
\int f_X(x)f_Y(z-x) \, dx
\]

\subsection*{3. Key Results from Lecture}

\begin{itemize}
\item Uniform + sine transformation $\rightarrow$ arcsine-type PDF
\item Sum of normals $\rightarrow$ Normal, variance adds
\item Sum of exponentials $\rightarrow$ Gamma-type PDF
\item Leibniz rule essential for differentiating integrals
\end{itemize}

\end{document}
